% Generated draft: requires mathematical and source review.
\documentclass[11pt]{article}
\usepackage[margin=25mm]{geometry}
\usepackage{fontspec}
\setmainfont{Latin Modern Roman}
\usepackage{amsmath,amssymb,mathtools,unicode-math}
\setmathfont{Latin Modern Math}
\usepackage{enumitem,needspace,xurl,hyperref}
\setlength{\parindent}{0pt}
\setlength{\parskip}{5pt}
\setlength{\emergencystretch}{3em}
\newcommand{\N}{\mathbb{N}}
\newcommand{\Z}{\mathbb{Z}}
\newcommand{\Q}{\mathbb{Q}}
\newcommand{\R}{\mathbb{R}}
\newcommand{\C}{\mathbb{C}}
\newcommand{\F}{\mathbb{F}}
\newcommand{\A}{\mathbb{A}}
\newcommand{\PP}{\mathbb{P}}
\newcommand{\E}{\mathbb{E}}
\newcommand{\eps}{\varepsilon}
\newcommand{\dd}{\,\mathrm{d}}
\begin{document}
\section*{UnsolvedMath ID: 4200013}
\textbf{Source problem code:} catalogue entry
\subsection{Definitions and mathematical statement}
The source describes an infinite-dimensional Calogero--Moser system of particles on the real line with inverse-square pair potential. For a countable-particle interpretation, formal equations with coupling $c$ are
\[
 \dot q_i=p_i,\qquad
 \dot p_i=2c\sum_{j\ne i}(q_i-q_j)^{-3},
\]
corresponding in finite dimension to $\tfrac12\sum_i p_i^2+c\sum_{i<j}(q_i-q_j)^{-2}$. The infinite index set, coupling convention/sign, phase space, summability or renormalization, and collision restrictions are not supplied. Thus these formulas explain the stated interaction but do not define a convergent infinite Hamiltonian or an already-existing flow. A kinetic/continuum limit is a different possible interpretation of the title's ``Vlasov'' term and is not identified with the particle system without additional data.

\textbf{Existence question (Oliver Knill).} Does the intended infinite-dimensional dynamics exist? A precise answer would have to specify the admissible states and the sense and time interval of a flow $\Phi^t$, including uniqueness if well-posedness is intended. No such initial-data class or local/global convention is given.

\textbf{Integrability question.} Is it integrable in the source's stated sense that \emph{every invariant measure} gives rise to almost-periodic dynamics? An invariant measure satisfies $(\Phi^t)_*\mu=\mu$. The source neither restricts to ergodic measures nor specifies whether measures are probabilities. For invariant probabilities, one possible measure-theoretic meaning is almost periodicity of the Koopman group: each orbit
\[
 \{h\circ\Phi^t:t\in\R\}\subset L^2(\mu)
\]
is relatively compact for every $h\in L^2(\mu)$. Almost periodicity of individual particle trajectories is another meaning and needs a state-space topology. These are not silently equated or substituted for the source's undefined notion.

The supplied summary claims substantial global existence and broad integrability results for continuum/kinetic Calogero--Moser dynamics, mentioning suitable spaces, mass constraints in a focusing case, conservation laws, and solitons. It gives no precise equation or citation matching the infinite inverse-square particle system. It separately says the every-invariant-measure almost-periodicity formulation is less addressed. Those existence/integrability notions are retained as distinct unverified claims, not used to settle the specific questions.
\subsection{Short English statement}
Construct the infinite inverse-square Calogero--Moser dynamics and ask whether every invariant measure yields almost-periodic dynamics. The particle or kinetic model, phase space, convergence, existence scope, and meaning of almost periodicity are unspecified. Reported continuum/PDE integrability is not identified with this exact measure-theoretic demand.
\subsection{Sources}
\begin{itemize}
\item Local UnsolvedMath registry, permanent ID 4200013, source code catalogue entry, ``Calogero--Moser--Vlasov''; proposer: Oliver Knill, proposed year 2000. Supplied status partially solved, classification PARTIAL-PROGRESS, updated 2026-08-06.
\item Knill, \emph{Open problems in Hamiltonian dynamics} (2000), item ``13) CALOGERO-MOSER-VLASOV,'' \url{https://people.math.harvard.edu/~knill/seminars/intr/}. Reported HTML extraction and access: 2026-07-29. The record says the source presents an open item and no later authoritative resolution check is recorded. Rights: NEEDS\_REVIEW; public access verified but redistribution permission not established.
\item The summary's continuum/kinetic well-posedness, focusing mass constraints, conserved-quantity, and soliton claims have no supplied papers or precise equations.
\end{itemize}
\end{document}
